% ============ 24 Poincaré 对偶 (原书 pp.144-146) ============
\chapter{Poincar\'e 对偶}\label{sec:24}
\index{Poincar\'e 对偶}
\index{迹映射}
\index{Gysin 映射}
\index{Verdier 对偶}
\index{定向层}
\index{投影公式}

\subsection*{拓扑空间的 Poincar\'e 对偶}

经典的 Poincar\'e 对偶定理(Greenberg, M., \emph{Lectures on Algebraic Topology}, Benjamin, 1967, 26.6)说，对定向连通 $m$ 维流形 $U$，有典范同构 $H^{r}_{c}(U, \mathbb{Z}/n\mathbb{Z}) \to H^{m-r}(U, \mathbb{Z}/n\mathbb{Z})$.利用 $H^{s}$ 与 $H_{s}$ 的对偶，可以把它改写为有限群的完美配对
\[
H^{r}_{c}(U, \mathbb{Z}/n\mathbb{Z}) \times H^{m-r}(U, \mathbb{Z}/n\mathbb{Z}) \to H^{m}_{c}(U, \mathbb{Z}/n\mathbb{Z}) \cong \mathbb{Z}/n\mathbb{Z}.
\]
更好的方法(例如 Iversen 1986, V.3)避免定向的选取.取而代之，引入``定向层'' $\mathcal{o}$，使得有典范同构 $H^{m}_{c}(U, \mathcal{o}) \cong \mathbb{Z}/n\mathbb{Z}$.Poincar\'e 对偶定理于是变成完美配对
\[
H^{r}_{c}(U, \mathcal{o}) \times H^{m-r}(U, \mathbb{Z}/n\mathbb{Z}) \to H^{m}_{c}(U, \mathcal{o}) \cong \mathbb{Z}/n\mathbb{Z}.
\]
流形可定向，当且仅当存在同构 $\mathbb{Z}/n\mathbb{Z} \to \mathcal{o}$，而这样的同构的选取就是 $U$ 的定向.

设 $C$ 表示 $\mathbb{R}$ 的代数闭包.给 $C$(看作实流形)一个定向与选取 $\sqrt{-1}$(从而对所有 $n$ 选取本原 $n$ 次单位根 $e^{2\pi\sqrt{-1}/n}$)是一回事.一旦选定 $C$ 的定向，就得到任何维数 1 的复流形的定向(共形映射保持定向)，事实上也得到任何维数的复流形的定向.于是，对复维数 $m$ 的连通复流形 $U$，一旦选定了本原 4 次单位根，经典的 Poincar\'e 对偶定理就取形式
\[
H^{r}_{c}(U, \mathbb{Z}/n\mathbb{Z}) \times H^{2m-r}(U, \mathbb{Z}/n\mathbb{Z}) \to H^{2m}_{c}(U, \mathbb{Z}/n\mathbb{Z}) \cong \mathbb{Z}/n\mathbb{Z}.
\]
更一般地，对任何有限维局部紧空间，有通常称为 Verdier 对偶的花哨对偶定理.

\subsection*{非奇异代数簇的 Poincar\'e 对偶}

设 $k$ 是代数闭域，$\Lambda = \mathbb{Z}/n\mathbb{Z}$，其中 $n$ 与 $k$ 的特征互素.对 $\Lambda$-模层 $F$，设 $\mathcal{L}F(m) = \mathcal{H}om(F, \Lambda(m))$.这里 $\Lambda(m) = \mu^{\otimes m}_{n}$.

设 $X$ 是维数 $d$ 的非奇异簇.对任何闭点 $P \in X$，Gysin 映射是同构
\[
H^{0}(P, \Lambda) \to H^{2d}_{P}(X, \Lambda(d)).
\]
有典范映射 $H^{2d}_{P}(X, \Lambda(d)) \to H^{2d}_{c}(X, \Lambda(d))$，我们设 $\mathrm{cl}(P)$ 是 $1$ 在这些映射的复合下的像.

\begin{theorem}\label{thm:24.1}
设 $X$ 是代数闭域 $k$ 上维数 $d$ 的非奇异簇.

(a) 有唯一映射 $\varphi(X) \colon H^{2d}_{c}(X, \Lambda(d)) \to \Lambda$，把 $X$ 上任何闭点 $P$ 的 $\mathrm{cl}(P)$ 送到 1；它是同构.$\varphi$ 称为迹映射.

(b) 对 $\Lambda$-模的任何局部常值层 $F$，有典范配对
\[
H^{r}_{c}(X, F) \times H^{2d-r}(X, \mathcal{L}F(d)) \to H^{2d}_{c}(X, \Lambda(d)) \cong \Lambda,
\]
它们是有限群的完美配对.
\end{theorem}

定义配对的一种方法是将 $H^{2d-r}(X, \mathcal{L}F(d))$ 等同于 $\operatorname{Ext}^{2d-r}_{X}(F, \Lambda(d))$，后者可以看作长度为 $2d - r$ 的扩张群，然后使用迭代余边缘映射(见定理~\ref{thm:14.20} 后的讨论).当 $X$ 拟射影时，这些群可以等同于 \v{C}ech 群，配对可以用通常的杯积公式定义.

\subsection*{Gysin 映射}

设 $\varphi \colon Y \to X$ 是代数闭域 $k$ 上光滑分离簇的紧合映射；设 $a = \dim X$，$d = \dim Y$, $e = d - a$.有限制映射
\[
\varphi^{*} \colon H^{2d-r}_{c}(X, \Lambda(d)) \to H^{2d-r}_{c}(Y, \Lambda(d)).
\]
由对偶，得到映射
\[
\varphi_{*} \colon H^{r}(Y, \Lambda) \to H^{r-2c}(X, \Lambda(-e)).
\]

\begin{remark}\label{rem:24.2}
(a) 映射 $\varphi_{*}$ 由方程唯一确定：
\[
\varphi_X(\varphi_{*}(y) \smile x) = \varphi_Y(y \smile \varphi^{*}(x)), \qquad x \in H^{2d-r}_{c}(X, \Lambda(d)), \ y \in H^{r}(Y, \Lambda).
\]
这就是定义.

(b) 若 $\varphi$ 是闭浸入 $Y \hookrightarrow Z$，则 $\varphi_{*}$ 是 \S 16 中定义的 Gysin 映射---这是对偶定理证明的结果.注意在这种情况下 $-e$ 是 $Y$ 在 $X$ 中的余维数.特别地，
\[
\varphi_{*}(1_Y) = \mathrm{cl}_X(Y),
\]
其中 $1_Y$ 是 $H^{0}(Y, \Lambda) = \Lambda$ 的单位元.

(c) 对复合映射
\[
\varphi_{1*} \circ \varphi_{2*} = (\varphi_1 \circ \varphi_2)_{*}
\]
这直接由定义推出，因为 $(\varphi_1 \circ \varphi_2)^{*} = \varphi^{*}_2 \circ \varphi^{*}_1$.

(d) 若 $Y$ 与 $X$ 完备，则
\[
\varphi_X(\varphi_{*}(y)) = \varphi_Y(y), \qquad y \in H^{2d}(Y, \Lambda(d)),
\]
因为
\[
\varphi_X(\varphi_{*}(y)) = \varphi_X(\varphi_{*}(y) \smile 1_X)
\]
\[
\xlongequal{(a)} \varphi_Y(y \smile \varphi^{*}(1_X)) = \varphi_Y(y \smile 1_Y) = \varphi_Y(y).
\]

(e) (投影公式) 若 $Y$ 与 $X$ 完备，则
\[
\varphi_{*}(y \smile \varphi^{*}(x)) = \varphi_{*}(y) \smile x, \qquad x \in H^{r}(X), \ y \in H^{s}(Y).
\]
为证明它，对 $x' \in H^{2d-r-s}(X)$ 应用 $\varphi_X$ 到 $\varphi_{*}(y \smile \varphi^{*}(x)) \smile x'$.

(f) 若 $\varphi \colon Y \to X$ 是次数为 $\delta$ 的有限映射，则 $\varphi_{*} \circ \varphi^{*} = \delta$.[习题] 由于 $\varphi^{*}$ 在 $H^{0}(X)$ 上恒等作用，$\varphi_{*}$ 在 $H^{2\dim X}(X)$ 上恒等作用.此外，$\varphi_{*}$ 在 $H^{2\dim X}(X)$ 上乘以 $\delta$，$\varphi^{*}$ 在 $H^{0}(X)$ 上乘以 $\delta$.

有时 $\varphi_{*}$ 也称为 Gysin 映射.
\end{remark}

\subsection*{对基变换定理的应用}

\begin{theorem}\label{thm:24.3}
设 $\varphi \colon Y \to X$ 是域 $k$(不一定是代数闭的)上簇的正则映射.设 $F$ 是 $Y$ 上的可构造层.则对所有 $r$，$R^{r}\varphi_{*}F$ 可构造，且除有限多个 $r$ 外为零.此外，$R^{r}\varphi_{*}$ 的形成与所有基变换 $T \to \operatorname{Specm} k$ 交换.
\end{theorem}

证明使用紧合基变换定理与 Poincar\'e 对偶.见 Deligne, P., \emph{Th\'eor\`emes de finitude en cohomologie $\ell$-adique}, SGA $4\frac{1}{2}$, pp.233--251.

\subsection*{Poincar\'e 对偶证明的梗概}

暂略.

\subsection*{Verdier 对偶}

暂略.关于花哨对偶定理的近期参考文献是：
\begin{quote}
Neeman, A., \emph{The Grothendieck duality theorem via Bousfield's techniques and Brown representability}, Jour. Amer. Math. Soc. 9 (1996), 205--236.
\end{quote}

\begin{quote}
\textbf{注记} 关于更一般的对偶定理，见 Riou 的讲义，\emph{Dualit\'e (d'apr\`es Ofer Gabber)}(\href{http://www.math.u-psud.fr/~riou/doc/dualite.pdf}{此处})与 Illusie 等人的讲义 p.126.
\end{quote}
